\section{The Dark Part of the Inventory}

A reader who follows astronomy at all has been waiting some pages to ask a question, and it deserves a straight answer rather than a silence. On the standard accounting, most of the matter in the universe is of a kind nobody has ever handled. It does not shine and does not block light; it announces itself only by its gravity --- in the speed of stars at the edges of galaxies, in the motions of galaxies within a cluster, in the bending of light around masses far greater than anything visible there. What does a universe of shadows and pushes say about that?

Less than you might hope, and I would rather set out exactly how much less than leave the impression of an answer.

Begin with what is \emph{not} a difficulty. This book is an account of how gravity works, not a census of what there is. In the shadow picture a body pulls its neighbour because it removes some of the rain that would otherwise have arrived from its direction --- so anything that stops massions gravitates, whether or not it has any dealings with light. Matter that blocks the flux and ignores our layer's photons is not forbidden here, nor even surprising; it is simply matter we cannot see. The toy takes the inventory from the astronomers and leaves it alone. That is not an evasion. It is the proper division of labour between a mechanism and a census.

Now the temptation, which is real, and which a sympathetic reader will feel before I do. This book's gravity has something the standard one does not: a natural distance, the length a massion travels before something stops it. Could the anomalies be \emph{that} --- gravity failing gently at great range --- rather than unseen matter? It would be a fine thing for the toy if they were. I have come to think they are not, and there are at least three good arguments against reaching for it, which are worth naming properly rather than gesturing at.

The first is about direction. If massions are absorbed along the way, then gravity at long range is \emph{screened} --- weaker than the inverse square, not stronger. What the rotation of galaxies demands is the opposite: more pull at the edges than the visible matter can supply. Absorption pushes the discrepancy the wrong way. It does not close the gap; it widens it.

The second is about what the anomaly keeps time with, and it is the argument I find hardest to answer. If the effect came from a mean free path it would set in at a certain \emph{distance} --- much the same distance in every galaxy, since the massions do not know whose galaxy they are crossing. It does not. Galaxies of very different sizes begin to misbehave at very nearly the same \emph{acceleration}, which is a pattern that a length cannot make. Whatever is happening is keyed to how hard gravity is pulling, not to how far one has gone.

The third is about separation. There are collisions of galaxy clusters in which the mass, as measured by the bending of light, has come apart from the visible matter and sits to either side of it. Any adjustment of the law of gravity is tied to where the visible matter \emph{is}; it has great difficulty producing gravity in a place where the visible matter is not.

So I do not think this book explains the dark matter away, and I would rather say that plainly than leave the possibility lying about as decoration.

But I want to be equally careful in the other direction, because \emph{this is not an alternative to dark matter} and \emph{this has nothing to do with dark matter} are different sentences, and I mean only the first. A tower of layers is a large and strange structure and I do not pretend to know everything that follows from it. Something a floor below us that engages the massions and not our light would be, in this book's own terms, precisely what the standard picture would record as mass that cannot be seen. Whether that thought leads anywhere I genuinely do not know --- I have not been able to make it do any work --- and I set it down as an unexplored hole in the wall rather than as an answer. Marked holes are more useful than papered ones.

One thing does come back from the exercise, and it belongs on the sixth wound's account. Gravity is observed to work across a galaxy, across a cluster of galaxies, and across the distances over which light is bent by masses lying most of the way to the edge of what we can see. If massions were being used up on the journey, gravity would fail somewhere along that range, and it does not. So the massion's mean free path must be at least as long as the greatest distance over which we can watch gravity operate --- which is a fourth demand on the stopping chance, arriving from the largest scales there are. It pulls the same way as the third rather than against it, so it does not tighten the pincer; what it does is stretch the range over which transparency must hold by a factor it is difficult to write down. A mechanical theory of gravity should be made to meet that, and this one is now on notice that it must.

Of the other darkness --- the accelerating expansion, and whatever is driving it --- this book has nothing whatever to say. It is an account of a mechanism, not a history of the universe, and it does not reach that far. I would rather be silent there than inventive.

